% =====================================================================
%  PETA KEWAJIBAN PER TAHAP  (HAPUS DARI NASKAH FINAL)
% =====================================================================
\halaman{PETA KEWAJIBAN PER TAHAP}

{\itshape\noindent Seminar proposal (TA 1): dilaksanakan sebelum penelitian;
yang dinilai adalah kelayakan rencana penelitian (Bab I sampai III). Sidang
akhir (TA 2): dilaksanakan setelah penelitian selesai; yang dinilai adalah
hasil, pembahasan, dan pemahaman mahasiswa.\par}

\vspace{0.2cm}
{\small\singlespacing\noindent \WAJIB: harus ada.
\textcolor{merahwajib}{\textbf{WAJIB bila berlaku}}: bila kondisinya terpenuhi.
\Opsional: boleh dikosongkan. \TidakPerlu: tidak diisi.\par}

% Makro sel berwarna
\newcommand{\cW}[1][WAJIB]{\cellcolor{bgwajib}\textcolor{merahwajib}{\textbf{#1}}}
\newcommand{\cO}{\cellcolor{bgops}\textcolor{kuningops}{\textbf{Opsional}}}
\newcommand{\cT}{\cellcolor{bgtidak}\textcolor{abuisian}{Tidak diperlukan}}

\vspace{0.3cm}
{\footnotesize
\renewcommand{\arraystretch}{1.15}
\setlength{\tabcolsep}{4pt}
\begin{longtable}{|>{\raggedright\arraybackslash}p{5.1cm}|
                   >{\centering\arraybackslash}p{2.45cm}|
                   >{\centering\arraybackslash}p{2.45cm}|}
\hline
\rowcolor{bgkepala}
\textbf{Bagian} & \textbf{Seminar proposal} & \textbf{Sidang akhir} \\ \hline
\endhead
Halaman judul, pengesahan, orisinalitas & \cW & \cW \\ \hline
Pernyataan penggunaan AI                & \cT & \cW \\ \hline
Kata pengantar                          & \cO & \cW \\ \hline
\textit{Abstract} dan Abstrak           & \cO & \cW \\ \hline
BAB I Pendahuluan                       & \cW & \cW[WAJIB (final)] \\ \hline
BAB II Tinjauan Pustaka                 & \cW & \cW[WAJIB (final)] \\ \hline
BAB III Metodologi                      & \cW & \cW[WAJIB (final)] \\ \hline
BAB IV Hasil dan Pembahasan             & \cT & \cW \\ \hline
BAB V Kesimpulan dan Saran              & \cT & \cW \\ \hline
Daftar pustaka                          & \cW & \cW \\ \hline
Lampiran 1 (Form Kendali Bimbingan)     & \cW & \cW \\ \hline
Lampiran 2 (berita acara seminar proposal) & \cT & \cW \\ \hline
Lampiran 3 sampai 5 (kode, peta kode, data) & \cO & \cW \\ \hline
Lampiran 6 sampai 8 (\textit{similarity}, log AI, riwayat versi) & \cT & \cW \\ \hline
Lampiran 9 (daftar riwayat hidup)       & \cT & \cO \\ \hline
\end{longtable}}

{\scriptsize\itshape\color{abuisian}\noindent Pembagian ini disusun dari
Pedoman TA Prodi IF dan rubrik TA non-konvensional; konfirmasikan ke
Koordinator TA bila ragu.\par}
\clearpage
